\subsection{SOP-Guided Alignment of LMM-as-Judge with Human Preference}
\label{sec:method_sop_alignment}

Developing a specialized, lightweight evaluation model for multi-subject personalized generation necessitates large-scale pairwise preference data. However, eliciting reliable preference labels in this domain is notoriously challenging. An evaluator must jointly verify the physical presence of multiple subjects, the fidelity of fine-grained identity and wardrobe attributes, and the correctness of complex cross-subject interactions. This multi-dimensional verification renders large-scale human annotation prohibitively expensive and highly susceptible to inter-annotator variance. 

To overcome this annotation bottleneck without compromising label quality, we propose an SOP-guided (Standard Operating Procedure) alignment framework. Specifically, we employ frontier Large Multimodal Models (LMMs) as scalable annotators to construct the training set, while reserving human annotators strictly for the evaluation and test sets. To bridge the policy gap between machine and human evaluators---a pervasive issue in the LMM-as-a-Judge paradigm---we enforce a unified SOP across both stages. By constraining both LMMs and humans to follow the exact same structured reasoning rubric, we ensure the lightweight evaluator is trained on machine-generated labels that faithfully reflect the human decision-making process.

\paragraph{SOP Formulation and Errors-First Scaffolding.}
Unlike generic preference prompting, which often collapses into uninterpretable holistic judgments (e.g., favoring global aesthetics over prompt faithfulness), our SOP decomposes the pairwise comparison into a systematic, three-stage verification process: 
\begin{enumerate}
    \item \emph{Existence}: verifying the physical presence of all referenced subjects.
    \item \emph{Appearance}: scrutinizing structural integrity and identity/wardrobe fidelity against the references.
    \item \emph{Interaction}: assessing the spatial arrangement and physical realism of the specified relations.
\end{enumerate}

Crucially, the SOP forces the judge to externalize its reasoning by generating a candidate-specific \emph{flaw log} before rendering a final binary preference. By explicitly penalizing missing entities, identity swapping, and interaction violations, this ``errors-first'' scaffolding acts as a regularizer, anchoring the LMM's latent preference policy to concrete compositional metrics rather than spurious visual correlations.

\paragraph{High-Confidence Label Aggregation via Cross-Model Consensus.}
To further distill high-fidelity training signals and mitigate model-specific biases, we implement a cross-model consensus mechanism. We independently query two distinct frontier models (\texttt{gemini-2.5-flash} and \texttt{gemini-3.1-flash-lite-preview}) using the identical SOP. We then construct the final training set by taking the intersection of their predictions---retaining only the pairwise comparisons where both models yield the exact same preference. While this intersection rule reduces the total yield, it acts as a rigorous precision filter, effectively pruning examples where preference is driven by idiosyncratic model noise or ambiguous image quality. The resulting high-agreement subset provides a robust, conservative proxy for human preference.

The exact prompt used to enforce this SOP in our pipeline is reproduced below. It explicitly separates independent candidate scoring from the final pairwise decision, defines strict edge cases, and mandates a JSON schema that forces the LMM to articulate flaws before concluding.

\begin{verbatim}
You are an expert judge for multi-subject personalized image generation.

You will receive:
1. The original generation prompt.
2. Reference images for each subject mentioned in the prompt.
3. Candidate image A.
4. Candidate image B.

Your job consists of two parts:
Part 1: Independent Evaluation
Evaluate candidate A and candidate B independently. For each image, strictly evaluate:
- Existence: Are ALL the required subjects in the references physically present?
  (Generic versions count as 1; completely missing counts as 0).
- Appearance: Are the subjects rendered WITHOUT structural deformations? Do the
  generated subjects strictly match the references across ALL features:
  Face/Skin tone, Hair style, Top clothing, AND Bottoms (pants)?
- Interaction: Does the generated image reflect the described actions, maintain
  correct proportions, and follow realistic spatial/physical laws?

Part 2: Overall Comparison
Compare Candidate A and B and decide the overall winner.

OUTPUT RULES:
- Use ONLY the provided prompt and images.
- All dimensional scores (`*_existence`, `*_appearance`, `*_interaction`) must be
  strictly binary: 1 (fully correct) or 0 (any issue exists).
- `winner` must be either "A" or "B" (Never "Tie").
- Return valid JSON ONLY. Absolutely NO markdown formatting.
- CRITICAL: To focus on errors, the JSON MUST begin with `a_flaw_log` and
  `b_flaw_log`. DO NOT list perfect subjects. ONLY list subjects with errors and
  explicitly state what is wrong (Missing, Wrong Face/Hair/Skin, Wrong Top,
  Wrong Pants, Action/Spatial/Physics Error). If an image is 100% perfect,
  write "None".

SCORING NOTES & SPECIFIC EDGE CASES:
- Appearance (Strict Identity & Wardrobe): You MUST check the whole person. If a
  subject has the correct top but wrong pants, wrong hair, or mismatched skin
  tone/face compared to the reference, score Appearance as 0 and log the
  specific flaw (e.g., "Woman red hoodie: Wrong face/skin color",
  "Man flannel shirt: Wrong pants color").
- Appearance (Mangled/Melted): Any structural artifacts (extra limbs, mangled
  faces) = Appearance 0.
- Existence vs Appearance: If requested "man denim jacket" is drawn as a generic
  man in a grey sweater, Existence is 1, Appearance is 0.
- Interaction (Partial): Evaluate interactions based ONLY on rendered subjects.
  If an interaction requires a missing subject, Interaction is 0.
- Interaction (Physics, Environment & Scale): Obvious physical violations
  (e.g., objects floating in mid-air without support) OR severe scale distortions
  between objects (e.g., abnormally giant bread, miniature people) = Interaction 0.
  HOWEVER, the presence of unprompted but logical supporting structures
  (e.g., tables, pedestals) used to hold objects is PERMITTED and should NOT be
  penalized.

EXPECTED JSON SCHEMA:
{
  "a_flaw_log": "Middle eastern man: Appearance (wears green sweater instead of
  beige shirt). Man flannel shirt: Appearance (wrong pants color). Woman red
  hoodie: Appearance (mismatched face/skin/hair). Interaction: Man is not
  occluded by the woman; The burger is floating in mid-air.",
  "b_flaw_log": "Middle eastern man: Missing. Man bomber jacket: Missing. DSLR
  camera: Missing. Interaction: Man is not walking towards/staring at the woman;
  Fails due to missing interaction targets.",
  "a_existence": 1,
  "a_appearance": 0,
  "a_interaction": 0,
  "b_existence": 0,
  "b_appearance": 0,
  "b_interaction": 0,
  "winner": "A"
}
\end{verbatim}
